% texte.tex — f171r, vv. 721–764 (conversion initiale des travaux antérieurs)
\vers[n=721, indent=2em, findent=1.2em]{}{}{ξὺν ἀθέοις ἔργοις}{Car toi, avec tes impiétés, nous te châtierons}
\vers[n=722, metre={tr⁴}]{}{}{καὶ γὰρ ἀνταμειψόμεσθά σ᾽ ὥσπερ εἰκὸς ἀντὶ τῶνδε.}{de ta conduite, comme de juste.}
\vers[n=723, indent=2em, findent=1.2em, metre={lyr.}]{}{}{τάχα δὲ μεταβαλοῦσ᾽ ἐπὶ κακὸν ἑτερότρο-}{Et puis la Fortune, changeant brusquement d'humeur aura vite fait de poursuivre un méchant homme}
\vers[n=725, indent=2em, findent=1.2em]{}{}{\alinea πον ἐπέχει τύχη.}{qui a pris le chemin du crime.}
\vers[n=726, metre={tr⁴}]{ΚΟ.}{Coryphée}{ἀλλὰ τάσδε μὲν λαβεῖν χρῆν σ᾽ ἐκφέρειν τε τῶν ξύλων,}{Mais prends ces femmes avec toi, et va chercher du bois.}
\vers[n=727]{}{}{καὶ καταίθειν τὸν πανοῦργον πυρπολεῖν θ᾽ ὅσον τάχος.}{Il faut le brûler vif, ce fourbe, l'anéantir au plus tôt par le feu.}
\vers[n=728]{ΜΙ.}{Mika}{ἴωμεν ἐπὶ τὰς κληματίδας ὦ Μανία.}{Allons chercher les fagots, Mania.}
\vers[n=729]{}{}{κἀγώ σ᾽ ἀποδείξω θυμάλωπα τήμερον.}{Je ferai de toi un tison aujourd'hui..}
\vers[n=730]{ΜΝ.}{Mnésiloque}{ὕφαπτε καὶ κάταιθε· σὺ δὲ τὸ Κρητικὸν}{Allumez, brûlez. Toi, ma petite, ôte vite ta crétique}
\vers[n=731]{}{}{ἀπόδυθι ταχέως· τοῦ θανάτου δ᾽ ὦ παιδίον}{et n'accuse de ta mort que ta mère.}
\vers[n=732]{}{}{μόνην γυναικῶν αἰτιῶ τὴν μητέρα.}{Qu'est ceci ? La fillette n'est plus qu'une outre}
\vers[n=733]{}{}{τουτὶ τί ἔστιν; ἀσκὸς ἐγένεθ᾽ ἡ κόρη}{pleine de vin, ... et avec cela chaussée de persiques.}
\vers[n=734]{}{}{οἴνου πλέως καὶ ταῦτα Περσικὰς ἔχων.}{Ô femmes archi-enfiévrées et archi-biberonnes,}
\vers[n=735]{}{}{ὦ θερμόταται γυναῖκες, ὦ ποτίσταται}{ingénieuses à trouver partout matière à boire !}
\vers[n=736]{}{}{κἀκ παντὸς ὑμεῖς μηχανώμεναι πιεῖν,}{Ô providence des débitants, et notre fléau à nous !}
\vers[n=737]{}{}{ὦ μέγα καπήλοις ἀγαθὸν ἡμῖν δ᾽ αὖ κακόν,}{Fléau aussi des ustensiles de ménage, et fléau de la trame !}
\vers[n=738]{}{}{κακὸν δὲ καὶ τοῖς σκευαρίοις καὶ τῇ κρόκῃ.}{(Les femmes reviennent avec des fagots.)}
\vers[n=739]{ΜΙ.}{Mika}{παράβαλλε πολλὰς κληματίδας ὦ Μανία.}{Entasse force fagots autour de lui, Mania.}
\vers[n=740]{ΜΝ.}{Mnésiloque}{παράβαλλε δῆτα· σὺ δ᾽ ἀπόκριναί μοι τοδί,}{Entasse, c'est bon. Mais toi, réponds à ceci.}
\vers[n=741]{}{}{τουτὶ τεκεῖν φῄς;}{Tu te dis mère de cet enfant ?}
\vers[suite]{ΜΙ.}{Mika}{καὶ δέκα μῆνας αὔτ᾽ ἐγὼ}{Dix mois durant}
\vers[n=742]{}{}{ἤνεγκον.}{je l'ai porté.}
\vers[suite]{ΜΝ.}{Mnésiloque}{ἤνεγκας σύ;}{Tu l'as porté, toi ?}
\vers[suite]{ΜΙ.}{Mika}{νὴ τὴν Ἄρτεμιν.}{Oui, par Artémis.}
\vers[n=743]{ΜΝ.}{Mnésiloque}{τρικότυλον; ἢ πῶς.}{(Il porte l'outre à ses lèvres.) Du trois cotyles ? Ou de quel prix ?}
\vers[suite]{ΜΙ.}{Mika}{εἰπέ μοι· τί μ᾽ ἠργάσω;}{Dis-moi : quel tour vas-tu me faire ?}
\vers[n=744]{}{}{ἀπέδυσας ὦναίσχυντέ μου τὸ παιδίον}{Impudent, tu as dévêtu mon enfant,}
\vers[n=745]{}{}{τυννοῦτον ὄν.}{qui est si menu.}
\vers[suite]{ΜΝ.}{Mnésiloque}{τυννοῦτο; μικρὸν νὴ Δία.}{Si menu ? (Il le regarde.) Il est petit, par Zeus.}
\vers[n=746]{}{}{πόσ᾽ ἔτη δὲ γέγονε; τρεῖς χοᾶς ἢ τέτταρας;}{Et quel âge a-t-il ? trois ou quatre choés ?}
\vers[n=747]{ΜΙ.}{Mika}{σχεδὸν τοσοῦτον χὤσον ἐκ Διονυσίων.}{À peu près, mais elle est ainsi depuis les dionysiaques.}
\vers[n=748]{}{}{ἀλλ᾽ ἀπόδος αὐτό.}{Allons ! rends-le-moi.}
\vers[suite]{ΜΝ.}{Mnésiloque}{μὰ τὸν Ἀπόλλω τουτονί.}{Non, par le glorieux Apollon.}
\vers[n=749]{ΜΙ.}{Mika}{ἐμπρήσομεν τοίνυν σε.}{Alors nous allons te brûler.}
\vers[suite]{ΜΝ.}{Mnésiloque}{πάνυ γ᾽ ἐμπίμπρατε·}{Parfaitement, brûlez.}
\vers[n=750]{}{}{αὕτη δ᾽ ἀποσφαγήσεται μάλ᾽ αὐτίκα.}{Mais j'égorge celle-ci à l'instant.}
\vers[n=751]{ΜΙ.}{Mika}{μὴ δῆθ᾽, ἱκετεύω σ᾽· ἀλλ᾽ ἔμ᾽ ὅ τι χρῄζεις ποίει}{Non, non, je t'en conjure. Fais à moi tout ce que tu voudras,}
\vers[n=752]{}{}{ὑπέρ γε τούτου.}{plutôt qu'à elle.}
\vers[suite]{ΜΝ.}{Mnésiloque}{φιλότεκνός τις εἶ φύσει.}{Tu as pour tes enfants une tendresse de naturel.}
\vers[n=753]{}{}{ἀλλ᾽ οὐδὲν ἧττον ἥδ᾽ ἀποσφαγήσεται}{Mais elle n'en sera pas moins égorgée.}
\vers[n=754]{ΜΙ.}{Mika}{οἴμοι τέκνον. δός τὸ σφαγεῖον Μανία,}{Ah! mon enfant ! Donne-moi une patère, Mania,}
\vers[n=755]{}{}{ἵν᾽ οὖν τό γ᾽ αἷμα τοῦ τέκνου τοὐμοῦ λάβω.}{que je recueille au moins le sang de ma fille..}
\vers[n=756]{ΜΝ.}{Mnésiloque}{ὕπεχ᾽ αὐτό, χαριοῦμαι γὰρ ἕν γε τοῦτό σοι.}{Tiens-la là-dessous. Je veux bien t'accorder cette seule grâce. (Il asperge Mika du vin de l'outre.)}
\vers[n=757]{ΜΙ.}{Mika}{κακῶς ἀπόλοι᾽, ὡς φθονερὸς εἶ καὶ δυσμενής.}{Puisses-tu périr de male mort! Que tu es jaloux et mauvais !}
\vers[n=758]{ΜΝ.}{Mnésiloque}{τουτὶ τὸ δέρμα τῆς ἱερείας γίγνεται.}{Cette peau revient à la prêtresse.}
\vers[n=759]{ΚΡ.}{Kritylla}{τί τῆς ἱερείας γίγνεται;}{Qu'est-ce qui revient à la prêtresse ?}
\vers[suite]{ΜΝ.}{Mnésiloque}{τουτί. λαβέ.}{Tiens, prends. (Il lui jette l'outre vide.)}
\vers[n=760]{ΚΡ.}{Kritylla}{ταλαντάτη Μίκα τίς ἐξεκόρησέ σε;}{Infortunée Mica, qui t'a souillée ?}
\vers[n=761]{}{}{τίς τὴν ἀγαπητὴν παῖδά σοὐξηράσατο;}{Qui t'a soustrait ta fille chérie ?}
\vers[n=762]{ΜΙ.}{Mika}{ὁ πανοῦργος οὗτος. ἀλλ᾽ ἐπειδήπερ πάρει,}{C'est ce fourbe. Mais puisque te voilà,}
\vers[n=763]{}{}{φύλαξον αὐτόν, ἵνα λαβοῦσα Κλεισθένη}{garde-le ; que j'aille, en compagnie de Clisthène,}
\vers[n=764]{}{}{τοῖσιν πρυτάνεσιν ἃ πεποίηχ᾽ οὗτος φράσω.}{raconter aux prytanes ce qu'il a fait.}
